\documentclass[12pt]{report}
\usepackage[a4paper,margin=1in]{geometry}
\usepackage{graphicx}
\usepackage{booktabs}
\usepackage{hyperref}
\usepackage{fancyhdr}
\usepackage{titlesec}
\usepackage{enumitem}

\pagestyle{fancy}
\fancyhf{}
\fancyhead[L]{UCS503 - Software Engineering}
\fancyhead[R]{\thepage}
\renewcommand{\headrulewidth}{0.4pt}

\titleformat{\chapter}[display]
  {\normalfont\huge\bfseries}{\chaptertitlename\ \thechapter}{20pt}{\Huge}

\begin{document}

% ---------------- TITLE PAGE ----------------
\begin{titlepage}
\centering
\vspace*{1cm}
{\large UCS503 - Software Engineering}\\[1.8cm]

{\Huge \bfseries Campus Marketplace}\\[0.4cm]
{\large A Verified Student Platform for Buying, Selling, Lending and Borrowing with AI-Driven Pricing and Fraud Detection}\\[2.2cm]

{\large Submitted by:}\\[0.4cm]
\begin{tabular}{@{}ll@{}}
1024031049 & \quad Shrasthi Singh Rathore \\
1024031021 & \quad Raksha Pandey \\
1024031047 & \quad Shweta Agarwal \\
\end{tabular}\\[0.6cm]
{\large Group: 3C73, BE Third Year, CSE}\\[2.2cm]

{\large Submitted to:}\\[0.4cm]
{\large Dr. Jhonsy Bansal}\\[2.2cm]

{\large Mid-Semester Evaluation}\\[0.4cm]
{\large\bfseries Computer Science and Engineering Department}\\
{\large Thapar Institute of Engineering and Technology, Patiala}
\vfill
\end{titlepage}

\tableofcontents

% ---------------- CHAPTER 1 ----------------
\chapter{Introduction}

\section{Background and Context}
College campuses generate a constant churn of used books, hostel items, cycles, electronics
and lecture notes that students want to sell, lend, buy or borrow, but there is no single
trusted place to do this. Existing general-purpose classifieds and social-media groups have
no way of verifying that a seller or buyer is actually a student of the same institute, no
structured record of listings, messages or transactions, and no way to judge whether a quoted
price is fair or whether a listing looks fraudulent. \textbf{Campus Marketplace} is a full-stack
web application built to solve this problem: a campus-exclusive marketplace where every user
is verified through their institutional email, and where an integrated machine-learning layer
assists with fair pricing and fraud screening.

\subsection{Problem Statement}
Students need a single, trustworthy platform, restricted to verified members of their own
institute, on which they can browse and publish listings for buying, selling, lending and
borrowing items; exchange messages with other users; track transactions and due dates for
borrowed items; and rely on the platform to flag unfair prices and suspicious listings, all
of this backed by persistent storage rather than data that disappears on refresh.

\section{Project Objectives}
\begin{enumerate}[label=\arabic*.]
    \item \textbf{Objective 1:} Build a fully functional MERN-stack backend (Express + MongoDB
    via Mongoose) so that listings, users, messages, transactions, reviews and moderation
    cases are persisted in a real database instead of hardcoded mock data.
    \item \textbf{Objective 2:} Restrict access to verified students by implementing
    institutional email verification during sign-up.
    \item \textbf{Objective 3:} Support the full buy/sell/lend/borrow lifecycle, including
    due-date tracking for items that are lent or borrowed rather than sold outright.
    \item \textbf{Objective 4:} Integrate machine-learning services for AI-based price
    recommendation and fraud risk scoring so that listings are priced fairly and suspicious
    activity is flagged automatically.
\end{enumerate}

% ---------------- CHAPTER 2 ----------------
\chapter{Methodology and Design}

\section{System Architecture}

\subsection{High-Level Design}
The application follows a modular, three-part architecture. A React (Vite) single-page
frontend handles routing, filtering and all user-facing screens (Buy/Borrow, Sell/Lend, My
Listings, Messages, Transactions, Profile and an Admin moderation view). An Express.js REST
API sits between the frontend and the database, exposing endpoints for listings, users,
messages, transactions, reviews and moderation, and enforcing business rules such as
seller identity and category metadata on new listings. MongoDB, accessed through Mongoose
models, is the persistent data store; during development, when a local MongoDB instance is
not available, the backend automatically falls back to an in-memory MongoDB instance so the
system remains fully runnable and demoable. Alongside this Node/Express core, independent
machine-learning microservices handle price recommendation and fraud scoring.

\subsection{Technical Stack}
\begin{itemize}
    \item \textbf{Frontend:} React (Vite), with Lucide React icons.
    \item \textbf{Backend:} Node.js, Express.js.
    \item \textbf{Database:} MongoDB (Mongoose ODM), with \texttt{mongodb-memory-server} as a
    development-time fallback, and MongoDB Atlas for hosted deployment.
    \item \textbf{Machine Learning Services:} Python, FastAPI, scikit-learn (Random
    Forest / XGBoost for price recommendation, Isolation Forest for fraud risk scoring).
    \item \textbf{Tools:} \texttt{cors} for cross-origin access, Vite dev-server proxying,
    npm scripts for running and seeding the backend.
\end{itemize}

\section{Implementation Details}

\subsection{Module 1: Core Functionality}
The core of this phase of the project was converting a frontend-only prototype into a fully
functional client-server application:
\begin{itemize}
    \item A dedicated \texttt{server/} package was created with Express and Mongoose, defining
    \texttt{User}, \texttt{Listing}, \texttt{Message}, \texttt{Transaction}, \texttt{Review}
    and \texttt{Moderation} models, and exposing REST routes under \texttt{/api/...}.
    \item Every major frontend page was wired to the API through a centralised
    \texttt{src/api.js} module, so that Buy/Borrow, Sell/Lend, My Listings, Messages,
    Transactions, Profile and Admin all read and write MongoDB instead of hardcoded arrays.
    \item Publishing a listing now creates a real document in MongoDB (\texttt{POST
    /api/listings}) with seller identity and category/fair-price metadata attached, and My
    Listings, Buy/Borrow and the Dashboard all refetch from the database, so new listings
    persist across page refreshes.
    \item ``Chat with seller'' and ``Buy/Borrow'' actions, which previously only showed a
    confirmation toast, now create real \texttt{Message} and \texttt{Transaction} documents,
    so conversation and purchase history accumulates with real use.
    \item Borrow/lend transactions carry a due date, and the application tracks and surfaces
    this due date so borrowed items can be returned on time and lenders can see when an item
    is expected back.
    \item CORS was enabled on the Express server and a Vite dev-server proxy was configured
    so the frontend (port 5173) can call \texttt{/api/...} and have it forwarded transparently
    to the API (port 5001), eliminating cross-origin and wrong-origin errors during
    development.
    \item Because Lucide icon components cannot be stored in MongoDB, listings store a string
    \texttt{iconKey} field (e.g. \texttt{"books"}, \texttt{"electronics"}); a
    \texttt{hydrateListing()} helper on the frontend maps this key back to the correct icon
    component before rendering.
    \item A seed script (\texttt{server/seed.js}) populates the database with realistic sample
    data (15 listings across books, cycles, hostel items, notes and electronics, along with
    users, message threads, transactions, reviews and moderation cases) automatically on
    first run when the listings collection is empty, so the UI is demo-ready even before real
    usage data accumulates.
\end{itemize}

\subsection{Module 2: Secondary Features}
On top of the core CRUD backend, the following supporting features have been completed:
\begin{itemize}
    \item \textbf{Institutional Email Verification:} Sign-up is restricted to users with a
    valid institutional email address, so that only verified students of the campus can list,
    message or transact on the platform.
    \item \textbf{AI Price Recommendation:} A FastAPI microservice hosts a regression model
    (Random Forest, with an XGBoost variant) trained to suggest a fair price band for a
    listing based on category, condition and comparable items, helping sellers price items
    reasonably and buyers judge whether an asking price is fair.
    \item \textbf{Fraud Detection:} A companion FastAPI microservice runs an Isolation Forest
    model to assign a fraud risk score to listings and transactions, flagging anomalous
    pricing or behaviour patterns for admin review through the moderation queue.
    \item \textbf{Borrow/Lend Due-Date Tracking:} Transactions created for lend/borrow
    listings (as opposed to outright sales) carry a due date, allowing both parties to track
    when a borrowed item is due for return.
\end{itemize}

% ---------------- CHAPTER 3 ----------------
\chapter{Results and Evaluation}

\section{Testing Methodology}
\begin{itemize}
    \item \textbf{Manual/Integration Testing:} Each API route was exercised end-to-end from
    the corresponding frontend page (e.g. publishing a listing from Sell/Lend and confirming
    it appears under My Listings and Buy/Borrow after a refresh).
    \item \textbf{Fallback Testing:} The in-memory MongoDB fallback was verified by running
    the server without a local \texttt{mongod} instance and confirming the API still started,
    connected and served data correctly.
    \item \textbf{Seed Verification:} The auto-seed logic was verified by starting against an
    empty database and confirming that listings, users, messages, transactions, reviews and
    moderation cases were created automatically.
\end{itemize}

\section{Performance Metrics}
The following was verified during a live run of the system:

\begin{table}[h]
\centering
\begin{tabular}{@{}ll@{}}
\toprule
\textbf{Item} & \textbf{Result} \\
\midrule
API server & Running on port 5001, connected in memory-DB mode \\
Seeded listings & 15 (books, cycles, hostel items, notes, electronics) \\
Seeded message threads & 3 \\
Seeded transactions & 3 \\
Seeded reviews & 3 \\
Seeded moderation cases & 3 \\
Frontend data source & Loaded live from MongoDB via React UI \\
\bottomrule
\end{tabular}
\caption{Verified system state after a clean start and auto-seed}
\end{table}

\section{Analysis}
All four project objectives set out in Section 1.2 were met for this phase: a working
persistent backend was built, institutional email verification restricts access to verified
students, the full buy/sell/lend/borrow lifecycle (including due-date tracking) is supported,
and both the AI price-recommendation and fraud-detection services are integrated. The
principal implementation challenges, and how each was resolved, are summarised below.

\begin{itemize}
    \item \textbf{No backend existed initially} -- resolved by building a dedicated
    Express + Mongoose \texttt{server/} package and wiring every frontend page to it through
    a single \texttt{api.js} module.
    \item \textbf{MongoDB was not installed on the development machine} -- resolved by adding
    an automatic \texttt{mongodb-memory-server} fallback, so the app runs identically whether
    or not a local database is available, while keeping a real \texttt{MONGODB\_URI} ready for
    production.
    \item \textbf{Frontend icon components could not be persisted} -- resolved by storing a
    string \texttt{iconKey} and hydrating it back to a Lucide icon on the client.
    \item \textbf{CORS/port mismatch between Vite and Express} -- resolved with \texttt{cors()}
    on the server and a Vite proxy for \texttt{/api/...}.
    \item \textbf{Empty database looked broken in demos} -- resolved with an auto-seed script
    that populates realistic sample data on first run.
    \item \textbf{Publishing/chat/buy actions did not persist} -- resolved by backing each
    action with a real MongoDB document (\texttt{Listing}, \texttt{Message},
    \texttt{Transaction} respectively).
\end{itemize}

% ---------------- CHAPTER 4 ----------------
\chapter{Conclusion}

\section{Summary of Work}
This phase converted Campus Marketplace from a frontend-only, mock-data prototype into a
fully functional MERN application with a persistent MongoDB backend, institutional email
verification, complete buy/sell/lend/borrow support with due-date tracking for borrowed
items, and two integrated machine-learning services for AI-based price recommendation and
fraud detection.

\section{Key Findings}
\begin{itemize}
    \item An in-memory MongoDB fallback is an effective way to keep a database-backed
    application fully runnable during development and demos when a local database
    installation is unavailable.
    \item Splitting concerns cleanly -- API layer, Mongoose models, seed script and frontend
    hydration -- made it straightforward to replace mock data with real persistence without a
    full frontend rewrite.
    \item Keeping the ML services (pricing and fraud detection) as independent FastAPI
    microservices, separate from the Node/Express core, kept the system modular and easier to
    test in isolation.
\end{itemize}

\section{Future Work and Recommendations}
The following enhancements are identified as valuable extensions beyond the current MVP:
\begin{enumerate}
    \item A dedicated mobile application.
    \item Computer-vision-based fraud detection (e.g. verifying listing photos).
    \item OCR-based invoice/receipt verification for transactions.
    \item An AI chat assistant to help users navigate listings.
    \item In-app escrow and payment handling.
    \item Demand forecasting to guide sellers on what to list and when.
    \item Opt-in geospatial discovery of nearby listings using MongoDB's 2dsphere indexing.
\end{enumerate}

\section{Final Remarks}
The hardest part of this phase was not the interface, which already existed, but making the
backend genuinely work end-to-end -- persisting data reliably even without a local MongoDB
installation, and turning previously stubbed actions (chat, buy, publish) into real,
persistent records. With the database, verification, due-date tracking, pricing model and
fraud-detection service all now in place, the platform is ready for the next phase of testing
and refinement.

\end{document}
